% Implementation — v0.2 prose.
% Target: ~3 pages. Engineering exposition: dual-implementation
% discipline, the corpus regression gate, the FP/FN ledger discipline,
% the reproducer commands.

\section{Implementation}
\label{sec:implementation}

\subsection{Dual implementation: Python and TypeScript}

% Source: docs/decisions/0023 (lockstep rule).
The agent-context AV story has two operational endpoints with
opposite runtime profiles. Server-side calls (registry publish,
the \verb|/check| public route, the validation pipeline) run on
Django and use Python detectors directly. Agent-side calls (the
MCP verification gate, the \verb|phub| CLI, browser tooling) run
on Node and use TypeScript detectors. Maintaining one canonical
implementation and a thin port at the other end produces drift
the moment a contributor edits one without the other; the
existing supply-chain literature treats this as a known
failure mode for security-critical code.

PrimitiveHub addresses it by treating cross-language verdict
parity as a CI-enforced contract, not a code-review claim. The
Architectural Decision Record that locks the contract
(ADR-0023~\cite{primitivehub:2026:adr-0023}) names two
requirements: (a) for any sample $s$ and detector $d$, the
verdict severity and the boolean \emph{passed} flag are
identical between the Python implementation in
\verb|apps/api/scanners/| and the TypeScript implementation in
\verb|packages/phub/src/scanners/|; (b) the per-detector
(TP, FP, TN, FN) tuples measured against the shared corpus must
agree byte-for-byte across engines. The contract holds at every
detector version we ship and is verified by
\verb|bench/scripts/parity_check.py| against the most recent
committed baseline JSONs.

The cost of the dual implementation is real: every detector
change is two PR edits and a parity-recheck. The benefit is
the same: the byte-identical-verdicts property is what lets
the MCP gate carry the same security guarantee as the server-side
validation pipeline, and it is what lets a reviewer reproduce
any verdict in either engine without converting between dialects
of the rules.

\subsection{The markdown-context analyzer}

% Source: docs/decisions/0035 (markdown context awareness).
Detectors that scan prose patterns share a markdown-context
analyzer that identifies three syntactic regions: YAML
frontmatter (delimited by \verb|---| at line starts), fenced
code blocks (opened by \verb|```| or \verb|~~~| at line starts
per the CommonMark spec~\cite{CommonMark:2024:spec}), and inline
code spans (delimited by backticks, with balanced run-length).
The analyzer walks the document once and produces three lists of
character offsets; per-detector calls into it are
constant-time per offset.

The analyzer's introduction (Round 1 of the elite-tier path)
was the v0.2.1 ratchet that closed three precision-defense FP
ledger entries simultaneously: detectors stopped firing on
documented teaching examples (fenced shell pipes), on inline-code
documentation of vendor commands (npx install in vendor docs),
and on YAML frontmatter referencing attack-shape keywords. ADR-0035
records the design and the trade-off: each detector opts into the
markdown-context skip by listing the detail tags that should
honor it, leaving the syntactic-attack-carrier families
(role-hijack via fenced \verb|system| markers, hidden
\verb|<!-- instructions -->| HTML comments) unprotected so they
still fire on their intended attack shape.

A second-order pass --- the semantic-context analyzer in
ADR-0036~\cite{primitivehub:2026:adr-0036} --- classifies the
\emph{prose} around a regex match into three categories:
negation (``never run \verb|curl|sh|''), mention (``Stripe's
typography''), and declaration (``CC-BY-4.0 (corpus content)''),
plus the v0.2.2 quoted-trigger predicate that recognises a
matched phrase wrapped in paired \verb|"...".| The semantic-context
analyzer is what lets detectors generalise beyond the syntactic
fence-skip: defender prose enumerating attacker phrases is now
distinguishable from attacker prose using them.

\subsection{The corpus regression gate}

% Source: docs/decisions/0037 (corpus as steering wheel);
% bench/baseline/FLOOR.md.
ADR-0037 reframes the corpus's role: not just a test set the
detectors are measured against, but the regression contract every
PR passes through. The mechanics:

\begin{itemize}
  \item The most recent committed baseline JSON under
        \verb|bench/baseline/| is the \emph{floor}. Per detector,
        a PR may not drop F\textsubscript{1} by more than the
        regression tolerance (0.02 absolute) against the floor
        on either engine without an explicit recall trade
        declared in the PR description.
  \item The shell script \verb|bench/scripts/gate.sh| is the
        single entry point. It discovers the latest baseline per
        language (lex-sort filename), runs the Py and TS
        measurement scripts with \verb|--gate --baseline <floor>|,
        and runs \verb|parity_check.py| on the resulting
        baselines.
  \item The script exits 0 iff no F\textsubscript{1} regression
        exceeds tolerance AND the cross-language tuple parity is
        intact.
  \item The GitHub Actions workflow
        \verb|.github/workflows/scanner-bench.yml| runs the same
        script on every PR touching detectors, audit pipelines,
        bench scripts, or workflow files. A regression fails the
        merge.
  \item A drift-gate step in the same workflow runs
        \verb|render-evaluation.py --check| to verify
        \verb|docs/evaluation.md| is in sync with the latest
        baselines, blocking merges where the public scoreboard
        would lie by being stale.
\end{itemize}

When the gate fires, the author has three options: close the
regression with a detector ratchet (lockstep TS + Py), re-label
the offending corpus sample with a FP-NNN or FN-NNN ledger
entry, or declare an explicit recall trade in the PR
description. Silent acceptance is not an option; the gate is
the steering wheel and the floor advances only by author intent.

\subsection{FP/FN methodology and ledgers}

% Source: docs/false-{positive,negative}-resolution.md.
The discipline that prevents tuning past optimal is the
append-only paired-ledger model PrimitiveHub adopts from
antivirus-industry practice~\cite{AntivirusTaxonomy:2010s:industry-practice}.
Every diagnosed finding (false positive or false negative) is
classified into exactly one of four resolution paths before any
detector change ships:

\begin{description}
  \item[Path A — edit the content.] The primitive author rewrites
    the prose so the detector no longer fires. Used when the
    finding is a real precision hazard caused by careless
    drafting on the primitive's side.
  \item[Path B — narrow allow entry plus ledger record.] A policy
    carve-out in \verb|.primitivehub/policy.yaml| documents that
    the finding is benign in this specific context. The allow
    entry MUST be accompanied by a numbered FP-NNN ledger entry
    that explains the carve-out's scope; an allow entry without
    a ledger entry is rejected by code review.
  \item[Path C — detector ratchet with corpus measurement.] The
    detector is structurally under-broad or over-broad. Lockstep
    edit to both implementations, accompanied by a corpus sample
    that locks the regression. The FN ledger entry shows the
    F\textsubscript{1} delta the fix produces.
  \item[Path D — escalate to a security reviewer.] The finding's
    classification is ambiguous; never silently allow.
\end{description}

The ledgers themselves
(\verb|docs/false-positive-resolution.md| and
\verb|docs/false-negative-resolution.md|) are append-only
markdown with numbered entries. As of this submission, the FP
ledger carries 14 entries (FP-001 through FP-014) and the FN
ledger carries 17 (FN-001 through FN-017). Each entry names
the corpus sample, classifies the path, records the resolution
vector, and (after closure) appends a re-review log showing the
F\textsubscript{1} delta and the lockstep-version bump. The two
ledgers together are the audit trail a peer reviewer can walk
to verify any specific claim in this paper's evaluation
section.

\subsection{Replicating the results}

% Source: bench/scripts/gate.sh; bench/cli/replay.py;
% docs/corpus-v0.1.md § "How to reproduce".
The corpus, detectors, measurement scripts, regression gate, and
per-sample replay tool are open-source under CC-BY-4.0 (corpus)
and Apache-2.0 (detectors and tooling). A reviewer reproduces
every number cited in Section~\ref{sec:evaluation} from a fresh
clone in under ten minutes:

\begin{lstlisting}[basicstyle=\ttfamily\small,frame=single,breaklines=true]
git clone https://github.com/DevXYZx/primitivehub.git
cd primitivehub
pnpm install --frozen-lockfile
pnpm --filter @primitivehub/phub build
(cd apps/api && uv sync --all-extras)
bench/scripts/gate.sh
python bench/scripts/render-evaluation.py --check
python bench/cli/replay.py --sample S-001
\end{lstlisting}

The gate emits per-detector P, R, and F\textsubscript{1} on both
engines and exits 0 iff the corpus regression contract holds.
The \verb|render-evaluation.py --check| step verifies that
\verb|docs/evaluation.md| (the public scoreboard) is in sync
with the latest baselines and exits 0 iff no drift is
detected.

The replay tool (\verb|bench/cli/replay.py|) lets a reviewer
re-run any individual sample through the Python detectors and
see the per-detector verdict with the outcome (TP, FN, FP, TN,
or ``cross-trip''-only) classification. With
\verb|--all --diff-only --include-holdout|, the tool processes
the full 109-sample corpus and reports only the samples whose
primary-detector outcome diverges from the manifest contract.
Holdout samples that surface ledger-documented gaps are tagged
\verb|[HOLDOUT]| and do not count toward the drift exit code,
keeping the tool useful as an open-loop replay surface
independent of the gate's blocking semantics.
